% Pursuing Stacks (cotes 134-2 to 134-8): archive cote 134-7, site page 63 (PDF page 64), typescript page 505.
% typescript. Draft compared: ps6.tex:3086-3138.
% AI-transcribed by Claude Opus 5.5 (claude-opus-5-5, adaptive thinking, effort high), Message Batches API, batch ps4 (msgbatch_012Aoc1NuskATPsyBeKSYK7J), 2026-11-12; not yet checked by eye.
\scanpage{134-7/63}
\mspage{505}

schematic homotopy types (over k=$\underline{Z}$) defined via\inserted[typed]{semisimplicial\hand{,}} unipotent bundles. Definit\unsure{e}
-ly, an understanding of schematic homotopy types will have to include the
(by now classical) Bousfield-Kan ideas, \inserted{and these} \deleted{which} are also relevant for my
reflections on``integration''and ``cointegration''\inserted[typed]{operations\hand{,}} (in connection with the
notion of a derivator (section 69)\hand{)}, called in their book ``homotopy \inserted{direct} \deleted{inverse}
limits'' and ``homotopy \inserted{inverse} \deleted{direct} limits'' (in the special case of the derivator
associated to ordinary homotopy types, if I got it right). It came as a
surprise that \inserted{in their book,} these operations are developped \xout{xx} mainly as technical tools
for developping \inserted[typed]{their} \xout{xxx} theory of k-completions, whereas in my own reflections
they appeared from the start as ``the'' main operations in homotopical as
well as homological algebra. There had been quite a similar surprise when
Tim Porter had sent me a reprint (in \over{\unsure{j}}{J}uly, just before I stopped \inserted{with my} \deleted{the} note\unsure{s}
for a month or so) of Don Anderson's beautiful paper ``Fibrations and \over{\unsure{g}}{G}eome-
tric \over{\unsure{r}}{R}ealizations'' (Bulletin of the Amer. Math. Soc. September 1978),
where a very general and (as I feel) quite basic existence theorem for
integration and cointegration operations\hand{,} of the type precisely I was af-
ter\hand{,} in the \xout{xxxxxf} set-up of closed model categories of Quillen\hand{,} is barely
alluded to at the end of the introduction, and \inserted[typed]{comes} \xout{xxxxxxxxxxappearsxasxa}
more or less as \inserted{just} a by\deleted{\unsure{e}}-product of \inserted[typed]{work done in view of a} \xout{\unsure{xxxexpositionxof}} a result on geometric
realizations which (to an outsider like me at any rate) looks highly te\hand{\unsure{c}}hni
-cal and not inspiring in\deleted{\unsure{t}} the least !

It is becoming clear that I cannot put off much longer getting acquain
-ted with the main ideas and results of Bousfield-Kan's book, which defini-
tely looks like one of the few basic texts \over{\unsure{i}}{o}n \xout{homotopyxtheory} foundational
matters in homotopy theory. Still, before doing some \inserted[typed]{basic} \xout{xxxx} reading, I
would like to write down the sporadic reflections on schematic homotopy
types I went into during the last \inserted{weeks} \deleted{month}, while they are still fresh in my
mind !

\margin{63}
